% ===========================================================================
% Math shortcuts, available in every formula on the site (blog posts and
% paper abstracts). Write them exactly as in a LaTeX preamble:
%
%   \newcommand{\R}{\mathbb{R}}
%   \newcommand{\norm}[1]{\left\lVert #1 \right\rVert}
%   \DeclareMathOperator*{\argmax}{argmax}
%
% \renewcommand, \def, \DeclareMathOperator, \DeclarePairedDelimiter,
% \definecolor, and \DeclareDocumentCommand (with m, o, and O{...}
% arguments) also work. Later definitions override earlier ones.
%
% Only math belongs here: theorem environments, \ref shortcuts, sectioning,
% and margin comments have no meaning inside a formula on the web. If a line
% can't be used, the page still works and the browser's JavaScript console
% says which line it was ("Math shortcut not loaded: ...").
% ===========================================================================

% ---- Compatibility, so preamble lines can be pasted in unchanged ----------
\newcommand{\ensuremath}[1]{#1}          % everything here is already math
\newcommand{\mathbbm}[1]{\mathbb{#1}}    % bbm package, for \mathbbm{1}
\newcommand{\nc}{\newcommand}
\newenvironment{displaymath}{}{}         % already display math inside $$...$$

% ---- Colours for coauthor notes -------------------------------------------
\definecolor{darkblue}{rgb}{0, 0, 0.55}
\definecolor{darkgreen}{rgb}{0, 0.39, 0}
\newcommand{\rh}[1]{{\color{darkblue}{#1}}}
\newcommand{\js}[1]{{\color{darkgreen}{#1}}}
\newcommand{\jl}[1]{{\color{purple}{#1}}}

% ---- Delimiters and operators ---------------------------------------------
\DeclarePairedDelimiter\ceil{\lceil}{\rceil}
\DeclarePairedDelimiter\floor{\lfloor}{\rfloor}
\DeclareMathOperator*{\argmax}{argmax}
\DeclareMathOperator*{\argmin}{argmin}
\DeclareMathOperator{\Hessian}{Hess}

% ---- Displays -------------------------------------------------------------
\newcommand{\bean}{\begin{eqnarray}}
\newcommand{\eean}{\end{eqnarray}}
\newcommand{\be}{\begin{displaymath}}
\newcommand{\ee}{\end{displaymath}}
\newcommand{\bea}{\begin{eqnarray*}}
\newcommand{\eea}{\end{eqnarray*}}

% ---- Divergences ----------------------------------------------------------
\newcommand{\kl}[2]{\mathrm{KL}(#1 \| #2)}
\newcommand*{\tv}[2]{\mathrm{TV}(#1, #2)}
\newcommand*{\hel}[2]{\mathrm{d_{H^2}}(#1, #2)}
\newcommand*{\chis}[2]{\chi^2(#1\| #2)}

% #2 set over the relation #3 (the optional first argument is ignored here)
\newcommand{\oset}[3][0ex]{\mathrel{\mathop{#3}\limits^{#2}}}

% ---- Brackets -------------------------------------------------------------
\newcommand{\brac}[1]{\left( #1 \right)}                  % brackets
\newcommand{\sbrac}[1]{\left\{ #1 \right\}}               % set
\newcommand{\abrac}[1]{\left\langle #1\right\rangle}      % angled

% ---- Probabilistic modifiers ----------------------------------------------
\newcommand{\lawof}{\mathscr{L}}
\newcommand{\lawequals}{\overset{\operatorname{law}}{=}}
\newcommand{\As}{\ensuremath{\text{a.s.}}}

% ---- Common distributions -------------------------------------------------
\newcommand{\Bernoulli}{\operatorname{Bern}}
\newcommand{\Geom}{\operatorname{Geom}}
\newcommand{\Exponential}{\operatorname{Exp}}
\newcommand{\Binomial}{\operatorname{Bin}}
\newcommand{\Poisson}{\operatorname{Pois}}
\newcommand{\Normal}{\cN}
\newcommand{\Unif}{\operatorname{Unif}}

% ---- Basic probability ----------------------------------------------------
\newcommand{\E}[1]{\mathbb{E}\left[#1\right]}                     % expectation
\newcommand{\sE}[2]{\mathbb{E}_{#2}\left[#1\right]}               % expectation w/ sub
\newcommand{\e}{{\mathbb E}}
\newcommand{\Var}[1]{\mathbf{Var}\left(#1\right)}                 % variance
\newcommand{\Cor}[2]{\mathbf{Cov}\left(#1,#2\right)}
\newcommand{\Varsub}[2]{\mathbf{Var}_{#2}\left(#1\right)}         % variance w/ sub
\newcommand{\Prob}[1]{\mathbb{P}\left(#1\right)}                  % probability
\newcommand{\Psub}[2]{{\mathbb P}_{#1}\left(#2\right)}            % probability subscript
\newcommand{\Psup}[2]{{\mathbf P}^{#1}\left\{#2\right\}}          % probability superscript
\newcommand{\I}[1]{{\mathbf 1}_{#1}}                              % indicator
\newcommand{\Cov}[2]{{\mathbf{Cov}}\left[#1, #2\right]}           % covariance
\newcommand{\Ent}[1]{\textbf{Ent}\left(#1\right)}
\newcommand{\Entsub}[2]{\mathbf{Ent}_{#2}\left(#1\right)}
\newcommand{\mP}[2]{{\mathbb P}_{#2}\left\{#1\right\}}
\newcommand{\mE}[2]{{\mathbb E}_{#2}\left[#1\right]}              % expectation (measure)
\newcommand{\mV}[2]{{\mathbf{Var}}_{#2}\left\{#1\right\}}         % variance (measure)
% \Cprob sizes the bracket on the term before conditioning; \probC after
\newcommand{\Cprob}[2]{\mathbb{P}\set{\left. #1 \; \right| \; #2}}
\newcommand{\probC}[2]{\mathbb{P}\left(#1 \; | \; #2\right)}
\newcommand{\CondE}[2]{\mathbb{E}\left[ #1 \; | \; #2 \right]}
\newcommand{\CondV}[2]{\mathbf{Var}\left( #1 \; | \; #2 \right)}
\newcommand{\phat}[1]{\ensuremath{\hat{\mathbf P}}\left\{#1\right\}}
\newcommand{\Ehat}[1]{\ensuremath{\hat{\mathbb E}}\left[#1\right]}
\newcommand{\ehat}{\ensuremath{\hat{\mathbf E}}}
\newcommand{\Esup}[2]{{\mathbb E^{#1}}\left\{#2\right\}}
\newcommand{\esup}[1]{{\mathbf E^{#1}}}
\newcommand{\Esub}[2]{{\mathbb E_{#1}}\left[#2\right]}
\newcommand{\esub}[1]{{\mathbf E_{#1}}}
\newcommand{\iid}{\text{i.i.d.}}
\newcommand{\as}{\text{a.s.}}

% ---- Convergence ----------------------------------------------------------
\DeclareDocumentCommand{\Pwto} {o} {
	\IfNoValueTF {#1}
	{\overset{n\to\infty}{\longrightarrow}}
	{ \xrightarrow[ #1 \to \infty]{\Pr }}
}
\DeclareDocumentCommand{\Prto} {o} {
	\IfNoValueTF {#1}
	{\overset{\Pr}{\longrightarrow}}
	{ \xrightarrow[ #1 \to \infty]{\Pr }}
}
\DeclareDocumentCommand{\Asto} {o} {
	\IfNoValueTF {#1}
	{\overset{\operatorname{a.s.}}{\longrightarrow}}
	{ \xrightarrow[ #1 \to \infty]{\operatorname{a.s.} }}
}
\DeclareDocumentCommand{\Mgfto} {o} {
	\IfNoValueTF {#1}
	{\overset{\operatorname{mgf}}{\longrightarrow}}
	{ \xrightarrow[ #1 \to \infty]{\operatorname{mgf} }}
}
\DeclareDocumentCommand{\Wkto} {o} {
	\IfNoValueTF {#1}
	{\overset{(d)}{\longrightarrow}}
	{ \xrightarrow[ #1 \to \infty]{(d) }}
}
\DeclareDocumentCommand \LPto { O{1} }
{\overset{\operatorname{L^{#1}}}{\longrightarrow}}

% restriction of the function #1 to #2
\newcommand{\restr}[2]{{\left.#1\vphantom{\big|}\right|_{#2}}}

% ---- Short names ----------------------------------------------------------
\nc{\on}{\operatorname}
\nc{\ch}{\mbox{ch}}
\nc{\Z}{{\mathbb Z}}
\nc{\C}{{\mathbb C}}
\nc{\prob}{\mathbb{P}}
\nc{\indic}{\mathbbm{1}}
\nc{\indistr}{\overset{d}{\longrightarrow}}
\nc{\inprob}{\overset{p}{\longrightarrow}}
\nc{\N}{{\mathbb N}}
\nc{\pone}{{\mathbb C}{\mathbb P}^1}
\nc{\pa}{\partial}
\nc{\F}{{\mathcal F}}
\nc{\arr}{\rightarrow}
\nc{\larr}{\longrightarrow}
\nc{\al}{\alpha}
\nc{\ri}{\rangle}
\nc{\lef}{\langle}
\nc{\W}{{\mathbb W}}
\nc{\la}{\lambda}
\nc{\lapl}{\mathcal{L}}
\nc{\ep}{\lambda}
\nc{\su}{\widehat{{\mathfrak s}{\mathfrak l}}_2}
\nc{\sw}{{\mathfrak s}{\mathfrak l}}
\nc{\g}{{\mathfrak g}}
\nc{\h}{{\mathfrak h}}
\nc{\n}{{\mathfrak n}}
\nc{\G}{\widehat{\g}}
\nc{\De}{\Delta}
\nc{\gt}{\widetilde{\g}}
\nc{\Ga}{\Gamma}
\nc{\one}{{\mathbf 1}}
\nc{\z}{{\mathfrak Z}}
\nc{\La}{\Lambda}
\nc{\wt}{\widetilde}
\nc{\wh}{\widehat}
\nc{\cri}{_{\kappa_c}}
\nc{\sun}{\widehat{\sw}_N}
\nc{\si}{\sigma}
\nc{\el}{\ell}
\nc{\om}{\omega}
\nc{\ol}{\overline}
\nc{\dzz}{\frac{dz}{z}}
\nc{\mc}{\mathcal}
\nc{\Cal}{\mathcal}
\nc{\bb}{{\mathfrak b}}
\nc{\ot}{\otimes}
\nc{\R}{{\mathbb R}}
\nc{\yy}{{\mc Y}}
\nc{\ga}{\gamma}
\nc{\us}{\underset}
\nc{\opl}{\oplus}
\nc{\Fq}{{\mathbb F}_q}
\nc{\Mq}{{\mathcal M}}
\nc{\Rep}{\on{Rep}}
\nc{\lan}{\langle}
\nc{\ran}{\rangle}
\newcommand{\ind}{\mathbbm{1}}

% ---- Calligraphic letters -------------------------------------------------
\newcommand\cA{\mathcal A}
\newcommand\cB{\mathcal B}
\newcommand\cC{\mathcal C}
\newcommand\cD{\mathcal D}
\newcommand\cE{\mathcal E}
\newcommand\cF{\mathcal F}
\newcommand\cG{\mathcal G}
\newcommand\cH{\mathcal H}
\newcommand\cI{\mathcal I}
\newcommand\cJ{\mathcal J}
\newcommand\cK{\mathcal K}
\newcommand\cL{{\mathcal L}}
\newcommand\cM{\mathcal M}
\newcommand\cN{\mathcal N}
\newcommand\cO{\mathcal O}
\newcommand\cP{\mathcal P}
\newcommand\cQ{\mathcal Q}
\newcommand\cR{{\mathcal R}}
\newcommand\cS{{\mathcal S}}
\newcommand\cT{{\mathcal T}}
\newcommand\cU{{\mathcal U}}
\newcommand\cV{\mathcal V}
\newcommand\cW{\mathcal W}
\newcommand\cX{{\mathcal X}}
\newcommand\cY{{\mathcal Y}}
\newcommand\cZ{{\mathcal Z}}

% ---- More short names -----------------------------------------------------
\nc{\D}{\mathcal D}
\nc{\Vect}{\on{Vect}}
\nc{\ghat}{\G}
\nc{\T}{\mc T}
\nc{\Tloc}{\T^\g_{\on{loc}}}
\nc{\vac}{|0\ran}
\nc{\Wick}{{\mb :}}
\nc{\mb}{\mathbf}
\nc{\delz}{\partial_z}
\nc{\K}{{\cali K}}
\nc{\cali}{\mathcal}
\nc{\li}{\mathfrak l}
\nc{\lt}{\widetilde{\li}}
\nc{\astar}{a^*}
\nc{\ka}{\kappa}

\nc{\OO}{{\mc O}}
\nc{\AutO}{\on{Aut}\OO}
\nc{\DerO}{\on{Der}\OO}
\nc{\DerpO}{\on{Der}_+\OO}
\nc{\Au}{{\mc A}ut}
\nc{\mf}{\mathfrak}
\nc{\V}{{\mc V}}
\nc{\hh}{\wh{\h}}

\nc{\pp}{{\mathfrak p}}
\nc{\mm}{{\mathfrak m}}
\nc{\rr}{{\mathfrak r}}
\nc{\gr}{\on{gr}}
\nc{\Spe}{\on{Spec}}
\nc{\rv}{\rho^\vee}
\nc{\can}{\on{can}}
\nc{\CC}{\on{Op}_G(D))}
\nc{\Op}{\on{Op}_G(D)}
\nc{\MOp}{\on{MOp}_G(D)}
\nc{\Db}{{\mathbb D}}
\nc{\ww}{w}

\nc{\oQl}{\ol{{\mathbb Q}}_\ell}
\nc{\oFq}{\ol{{\mathbb F}}_q}
\nc{\Q}{{\mathbb Q}}
\nc{\Ql}{{\mathbb Q}_\ell}
\nc{\bs}{\backslash}
\nc{\AD}{{\mathbb A}}
\nc{\M}{{\mc M}}
\nc{\Bun}{\on{Bun}}
\nc{\hl}{h^{\leftarrow}}
\nc{\hr}{h^{\rightarrow}}
\nc{\supp}{\on{supp}}
\nc{\He}{\on{H}}
\nc{\Aut}{\on{Aut}}
\nc{\Ll}{{\mc L}}
\nc{\Coh}{{{\mathcal C}oh}}
\nc{\ovc}{\overset{\circ}}
\nc{\Hav}{\on{H}}
\nc{\Mod}{\on{Mod}}
\nc{\kk}{{\mathfrak k}}
\nc{\vf}{\varphi}
\nc{\Gr}{\on{Gr}}
\nc{\gen}{\on{gen}}
\nc{\IC}{\on{IC}}
\nc{\Jac}{\on{Jac}}

% ---- Used on this site before the list above was added --------------------
\newcommand{\eps}{\varepsilon}
\newcommand{\norm}[1]{\left\lVert #1 \right\rVert}
\newcommand{\abs}[1]{\left\lvert #1 \right\rvert}
\newcommand{\ip}[2]{\left\langle #1, #2 \right\rangle}
